\chapter{Kommentierte Quellen}

Achtundzwanzig Werke, eingeordnet: Zu jedem steht, wofür man es aufschlägt und was es leistet, das andere nicht leisten.

\lit{1}{Arrow, K.: \textit{Social Choice and Individual Values}. 2.\ Aufl., Yale 1963.}{Der Unmöglichkeitssatz aus Kapitel~13 im Original --- ein schmaler Band, aus einer Doktorarbeit hervorgegangen, der die Wohlfahrtsökonomie umgepflügt hat. Der Beweis ist kurz und schwer; die Einleitung erklärt, warum die Bedingungen alle plausibel sind.}

\lit{2}{Bellman, R.: \textit{Dynamic Programming}. Princeton 1957.}{Das Buch, das die Rückwärtsrechnung aus Kapitel~7 begründete. \textbf{Bellman wählte den Namen nach eigener Auskunft, weil \enquote{Programmierung} damals unverdächtig klang} und er einem skeptischen Geldgeber nicht mit \enquote{Mathematik} kommen wollte.}

\lit{3}{Bernoulli, D.: \textit{Specimen theoriae novae de mensura sortis}. Petersburg 1738.}{Die Geburtsurkunde des Nutzenbegriffs, auf Latein und knapp zwanzig Seiten. Enthält das Petersburger Spiel, den Logarithmus als Nutzenfunktion und --- bemerkenswert früh --- die Einsicht, dass der Wert eines Gewinns vom Vermögen des Gewinners abhängt.}

\lit{4}{Briggs, A., Claxton, K., Sculpher, M.: \textit{Decision Modelling for Health Economic Evaluation}. Oxford 2006.}{\textbf{Das praktische Handbuch zu den Kapiteln~5, 6 und 10.} Entscheidungsbäume, Markow-Modelle, probabilistische Analyse, Akzeptanzkurven --- mit Tabellenblättern zum Nachbauen. Wer selbst ein Modell rechnen will, beginnt hier.}

\lit{5}{Claxton, K., Sculpher, M.: \textit{Using value of information analysis to prioritise health research}. PharmacoEconomics 24, 2006.}{Der Wert der Information aus Kapitel~3, angewandt auf die Forschungsförderung. Zeigt, wie man aus dem EVPI je Patient und der Zahl der Betroffenen eine Obergrenze für sinnvolle Studienkosten gewinnt.}

\lit{6}{Drummond, M. u.\,a.: \textit{Methods for the Economic Evaluation of Health Care Programmes}. 4.\ Aufl., Oxford 2015.}{Das Standardwerk der gesundheitsökonomischen Bewertung --- Kostenarten, Perspektiven, Diskontierung, Kosten-Wirksamkeit. Behandelt sorgfältig, was Kapitel~10 nur andeuten konnte: dass der Blickwinkel die Rechnung bestimmt.}

\lit{7}{Ellsberg, D.: \textit{Risk, Ambiguity, and the Savage Axioms}. Q.\ J.\ Econ.\ 75, 1961.}{Das Urnenbeispiel aus Kapitel~11 --- der Nachweis, dass Menschen bekannte Wahrscheinlichkeiten unbekannten vorziehen, auch wenn keine Zuordnung das erklären kann. \textbf{Derselbe Ellsberg veröffentlichte zehn Jahre später die Pentagon-Papiere.}}

\lit{8}{Elwyn, G. u.\,a.: \textit{Shared Decision Making: A Model for Clinical Practice}. J.\ Gen.\ Intern.\ Med.\ 27, 2012.}{Die gemeinsame Entscheidungsfindung in drei Schritten --- Wahlmöglichkeit benennen, Möglichkeiten erläutern, Vorlieben erkunden. \textbf{Die praktische Umsetzung dessen, was Kapitel~1 als Trennung von Zahlen und Werten begründet.}}

\lit{9}{Gaertner, W.: \textit{A Primer in Social Choice Theory}. 2.\ Aufl., Oxford 2009.}{Die Fundstelle für den Beweis des Satzes von Arrow und für die Wege daran vorbei --- eingeschränkte Vorliebenbereiche, Nutzenvergleiche, Mehrheitsregeln. Knapp und präzise; die Vertiefung zu Kapitel~13.}

\lit{10}{Gigerenzer, G.: \textit{Risiko. Wie man die richtigen Entscheidungen trifft}. Bertelsmann 2013.}{\textbf{Auf Deutsch und für jeden lesbar.} Warum natürliche Häufigkeiten verständlicher sind als Prozentzahlen, wie Ärzte und Patienten Testergebnisse missdeuten, und warum einfache Faustregeln oft genügen. Das Gegenstück zu Kapitel~12 --- und dessen praktische Nutzanwendung.}

\lit{11}{Gold, M., Siegel, J., Russell, L., Weinstein, M.: \textit{Cost-Effectiveness in Health and Medicine}. Oxford 1996.}{Der Bericht der amerikanischen Sachverständigenkommission, der die Regeln festlegte: gesellschaftlicher Blickwinkel, drei Prozent Diskontierung, QALY als Maß. \textbf{Viele Konventionen der Kapitel~6 und 10 stammen aus diesem Buch} --- und sind Konventionen geblieben, keine Ergebnisse.}

\lit{12}{Hunink, M. u.\,a.: \textit{Decision Making in Health and Medicine}. 2.\ Aufl., Cambridge 2014.}{\textbf{Das Lehrbuch der klinischen Entscheidungsanalyse und die Hauptquelle dieses Bandes.} Schwellenwerte, Bäume, Markow-Modelle, Informationswert, Nutzenerhebung --- alles am klinischen Beispiel, mit Aufgaben. Wer nur ein Buch zu diesem Band liest, lese dieses.}

\lit{13}{Kahneman, D.: \textit{Schnelles Denken, langsames Denken}. Siedler 2012.}{\textbf{Auf Deutsch.} Der Rückblick auf vierzig Jahre Forschung über tatsächliche Entscheidungen --- Bezugspunkte, Verlustscheu, Rahmungseffekte. Erzählend geschrieben; die Belege stehen in den Originalarbeiten, nicht hier.}

\lit{14}{Kahneman, D., Tversky, A.: \textit{Prospect Theory: An Analysis of Decision under Risk}. Econometrica 47, 1979.}{Die Originalarbeit zu Kapitel~12 und die meistzitierte der Ökonometrie überhaupt. Enthält die Wertfunktion mit Knick am Bezugspunkt und die Gewichtung kleiner Wahrscheinlichkeiten --- beides unmittelbar an Befragungen abgelesen.}

\lit{15}{Keeney, R., Raiffa, H.: \textit{Decisions with Multiple Objectives}. Cambridge 1993.}{Das Werk zu Kapitel~10: Wie man mehrere Ziele in eine Bewertung bringt, wann das überhaupt zulässig ist und welche Unabhängigkeitsbedingungen dafür nötig sind. Gründlich und anspruchsvoll.}

\lit{16}{Lattimore, T., Szepesvári, C.: \textit{Bandit Algorithms}. Cambridge 2020.}{Die vollständige Theorie zu Kapitel~9 --- obere Vertrauensschranken, Thompson-Ziehung, untere Schranken, Beweise. Die Fundstelle für die Schranke von Lai und Robbins. Frei verfügbar; mathematisch fordernd.}

\lit{17}{von Neumann, J., Morgenstern, O.: \textit{Theory of Games and Economic Behavior}. 2.\ Aufl., Princeton 1947.}{Satz 2.1 steht im Anhang der zweiten Auflage. \textbf{Das Hauptwerk gehört zu Band 21}; der Anhang begründete ein eigenes Fach. Die Axiome sind dort anders nummeriert und schwerer zu lesen als in jeder späteren Darstellung.}

\lit{18}{NICE: \textit{Health Technology Evaluations: The Manual}. London 2022.}{Das Regelwerk der britischen Bewertungsbehörde --- Schwellenwert, Diskontierung, Unsicherheitsanalyse, Sonderregeln für lebensverlängernde Behandlungen am Lebensende. \textbf{Das lesenswerteste Beispiel dafür, wie aus den Rechnungen dieses Bandes verbindliche Verfahren werden} --- samt der Ausnahmen, die man einbauen muss.}

\lit{19}{Pauker, S., Kassirer, J.: \textit{Therapeutic decision making: a cost-benefit analysis}. N.\ Engl.\ J.\ Med.\ 293, 1975.}{Die Behandlungsschwelle aus Kapitel~1, im Original und für Kliniker geschrieben. Sechs Seiten, ohne Formalismus, mit Rechenbeispielen --- der Beginn der klinischen Entscheidungsanalyse.}

\lit{20}{Pauker, S., Kassirer, J.: \textit{The threshold approach to clinical decision making}. N.\ Engl.\ J.\ Med.\ 302, 1980.}{Die Erweiterung auf Tests: zwei Schwellen, ein Testbereich. \textbf{Satz 4.1 in seiner ursprünglichen Fassung} --- und noch immer die klarste Darstellung des Gedankens, dass ein Test außerhalb dieses Bereichs nichts ändert.}

\lit{21}{Puterman, M.: \textit{Markov Decision Processes}. Wiley 1994.}{Das Standardwerk zu Kapitel~8 --- Wertiteration, Strategieiteration, Konvergenzbeweise, durchschnittliche gegen abgewertete Belohnung. Die Fundstelle für die ausgelassene Konvergenz; trocken, vollständig, unentbehrlich.}

\lit{22}{Raiffa, H.: \textit{Decision Analysis: Introductory Lectures on Choices under Uncertainty}. Addison-Wesley 1968.}{Die Vorlesung, die das Fach lehrbar machte --- Entscheidungsbäume, Informationswert, Nutzenerhebung, alles anschaulich und mit Sorgfalt in der Darstellung. Sechzig Jahre alt und didaktisch bis heute unübertroffen.}

\lit{23}{Savage, L.: \textit{The Foundations of Statistics}. Wiley 1954.}{Die Begründung subjektiver Wahrscheinlichkeiten aus Vorlieben --- was von Neumann für den Nutzen leistete, leistet Savage für die Wahrscheinlichkeit. \textbf{Damit ist auch die zweite Hälfte von Definition 1.1 axiomatisch gedeckt}, und der Bayessche Zugang wird zur Folge von Vernunftregeln.}

\lit{24}{Schöffski, O., Graf von der Schulenburg, J.-M.: \textit{Gesundheitsökonomische Evaluationen}. 4.\ Aufl., Springer 2012.}{\textbf{Auf Deutsch und mit Blick auf das deutsche System.} Kostenrechnung, Nutzenmessung, Effizienzgrenze statt Schwellenwert --- die deutsche Besonderheit, die in den englischen Werken fehlt. Die Ergänzung zu [6] und [18] für hiesige Verhältnisse.}

\lit{25}{Sonnenberg, F., Beck, J.: \textit{Markov models in medical decision making}. Med.\ Decis.\ Making 13, 1993.}{Die praktische Einführung in Kapitel~6 --- Zustände, Zyklen, Halbzykluskorrektur, Fallstricke. Zwölf Seiten, die jeder liest, der ein solches Modell baut, und die man in zwei Stunden durcharbeiten kann.}

\lit{26}{Sutton, R., Barto, A.: \textit{Reinforcement Learning: An Introduction}. 2.\ Aufl., MIT Press 2018.}{Die Verbindung von Kapitel~8 und 9 zum maschinellen Lernen aus Band 22: Wie man optimale Strategien lernt, statt sie zu berechnen. Frei verfügbar, anschaulich, mit dem Bandit-Problem als erstem Kapitel.}

\lit{27}{Torrance, G.: \textit{Measurement of health state utilities for economic appraisal}. J.\ Health Econ.\ 5, 1986.}{Die Übersicht über Standardlotterie, Zeitabwägung und Skala aus Kapitel~2 --- mit den Gründen, warum die drei Verfahren systematisch verschiedene Werte liefern. Immer noch die beste Darstellung dieser unangenehmen Tatsache.}

\lit{28}{Wegwarth, O., Gigerenzer, G.: \textit{Statistisches Denken in der Medizin}. Dtsch.\ Ärztebl.\ 110, 2013.}{\textbf{Auf Deutsch, vier Seiten, frei verfügbar.} Wie man Nutzen und Schaden gegenüber Patienten so darstellt, dass sie verstanden werden --- absolute statt relative Zahlen, natürliche Häufigkeiten, Zeitbezug. \textbf{Die kürzeste sinnvolle Antwort auf Kapitel~12.}}

\begin{merksatz}
\boxtitel{Wenn Sie nur drei Werke lesen}
\begin{center}
\renewcommand{\arraystretch}{1.6}
\begin{tabular}{@{}P{1.0cm}P{3.4cm}P{3.8cm}@{}}
\toprule
\textbf{[12]} & \textbf{Hunink u.\,a.} & \textbf{Für alles}\\
\textbf{[20]} & \textbf{Pauker, Kassirer} & Für die Schwellen\\
\textbf{[10]} & \textbf{Gigerenzer} & \textbf{Für das Gespräch}\\
\bottomrule
\end{tabular}
\end{center}

\begin{center}
\begin{minipage}{0.88\textwidth}
\kursiv{Hunink und Mitarbeiter lehren das Handwerk; Pauker und Kassirer liefern auf sechs Seiten den Gedanken, der die meisten unnötigen Untersuchungen erspart. \textbf{Und Gigerenzer zeigt, wie man das Ergebnis so mitteilt, dass es ankommt} --- was nach Kapitel~12 kein Beiwerk ist, sondern Teil der Entscheidung.}
\end{minipage}
\end{center}

\begin{center}
\begin{minipage}{0.88\textwidth}
\kursiv{Nach Interessenlage: [4] Briggs u.\,a.\ zum Modellbau, [16] Lattimore und Szepesvári für die Theorie des Erkundens, [9] Gaertner für die Gruppenentscheidung, [18] NICE für die Praxis der Bewertung. Auf Deutsch: [24] Schöffski und Graf von der Schulenburg für die Gesundheitsökonomie, [13] Kahneman und [28] Wegwarth und Gigerenzer für das tatsächliche Entscheiden.}
\end{minipage}
\end{center}
\end{merksatz}

\bigskip
\begin{flushright}
\itshape\color{grau}
Die Zahlen sagen, was wahrscheinlich ist.\\
Die Werte sagen, was zählt.\\
Erst beides zusammen sagt,\\
was zu tun ist.
\end{flushright}
